% LaTeX report: First (injection) dataset generation for GNN2
% Device-level shares are derived from nominal loads in ieee34Mod1_with_loadshape.dss (summed and normalized per device).

\documentclass[11pt]{article}
\usepackage[utf8]{inputenc}
\usepackage{amsmath,amssymb,bm}
\usepackage{geometry}
\usepackage{booktabs}
\usepackage{enumitem}
\geometry{margin=1in}

\title{Injection Dataset Generation (GNN2 First Dataset)}
\author{GNN2 Project}
\date{}

\begin{document}
\maketitle

\section{Overview}

The first dataset is the \textbf{injection dataset}: each sample is one AC power-flow snapshot with node-level \textbf{inputs} $(P_{\mathrm{inj},i}, Q_{\mathrm{inj},i})$ (kW, kvar) and \textbf{targets} voltage magnitude $|V_i|$ (pu) and angle $\angle V_i$ (deg) at each bus-phase node $i$. The graph is the static phase-level network (lines and transformers); features are the net injections at each node. At the source bus, the injection is the \textbf{power from the upstream grid} (per phase); at all other nodes it is net local (PV minus load, plus capacitors for $Q$). Device-level load/PV shares are derived from the nominal loads in the existing OpenDSS file \texttt{ieee34Mod1\_with\_loadshape.dss}: nominal P (and Q) per load device were summed to get system totals, then each device's share was set to its nominal over the total, giving \texttt{DEVICE\_P\_SHARE} and \texttt{DEVICE\_Q\_SHARE}.

\section{Scenario sampling}

Each \textbf{scenario} corresponds to one day. The following scalars are drawn \textbf{once per scenario} from uniform scaling factors applied to a baseline:
\begin{equation}
  \theta_s \sim \mathrm{Uniform}(\underline{r}_\theta,\, \overline{r}_\theta), \qquad
  P_{\mathrm{load}} = P_{\mathrm{load}}^{(0)} \cdot \theta_{P_{\mathrm{load}}}, \quad
  Q_{\mathrm{load}} = Q_{\mathrm{load}}^{(0)} \cdot \theta_{Q_{\mathrm{load}}}, \quad
  P_{\mathrm{pv}} = P_{\mathrm{pv}}^{(0)} \cdot \theta_{P_{\mathrm{pv}}}, \quad
  \sigma_{\mathrm{load}} = \sigma_{\mathrm{load}}^{(0)} \cdot \theta_{\sigma_{\mathrm{load}}}, \quad
  \sigma_{\mathrm{pv}} = \sigma_{\mathrm{pv}}^{(0)} \cdot \theta_{\sigma_{\mathrm{pv}}},
\end{equation}
with baseline $(P_{\mathrm{load}}^{(0)}, Q_{\mathrm{load}}^{(0)}, P_{\mathrm{pv}}^{(0)}, \sigma_{\mathrm{load}}^{(0)}, \sigma_{\mathrm{pv}}^{(0)})$ and ranges $(\underline{r}, \overline{r})$ per parameter (e.g.\ $P_{\mathrm{load}}$ in $[0.8, 1.2]$, $P_{\mathrm{pv}}$ in $[0.6, 1.2]$). $\sigma_{\mathrm{load}}$ and $\sigma_{\mathrm{pv}}$ are clipped to $[0, 0.5]$ and control per-device multiplicative noise at each snapshot (see below).

\section{Daily profiles and time resolution}

\begin{itemize}[noitemsep]
  \item \textbf{Time grid}: $T = 288$ steps per day (5-minute resolution, $\Delta t = 5$ min).
  \item \textbf{Load profile}: $m_L[t] \in \mathbb{R}^T$ from the DSS loadshape (e.g.\ \texttt{5minDayShape}, CSV).
  \item \textbf{PV profile}: $m_{\mathrm{PV}}[t] \in \mathbb{R}^T$ from the DSS irradiance shape (e.g.\ \texttt{IrradShape}, CSV).
  \item \textbf{Net profile} (for time stratification): $p_{\mathrm{net}}[t] = P_{\mathrm{load}}\, m_L[t] - P_{\mathrm{pv}}\, m_{\mathrm{PV}}[t]$.
\end{itemize}

At each integer time index $t \in \{0,\ldots,T-1\}$, system totals are:
\begin{equation}
  P_{\mathrm{load}}(t) = P_{\mathrm{load}}\, m_L[t], \qquad
  Q_{\mathrm{load}}(t) = Q_{\mathrm{load}}\, m_L[t], \qquad
  P_{\mathrm{pv}}(t) = P_{\mathrm{pv}}\, m_{\mathrm{PV}}[t].
\end{equation}

\section{Time-step selection (stratified)}

Not all $T$ steps are used per scenario. A total of $K_{\mathrm{total}}$ time indices are selected and split across three profiles (load, PV, net) so that low, medium, and high values of each profile are represented:
\begin{itemize}[noitemsep]
  \item $K_{\mathrm{load}}$, $K_{\mathrm{pv}}$, $K_{\mathrm{net}}$ with $K_{\mathrm{load}} + K_{\mathrm{pv}} + K_{\mathrm{net}} = K_{\mathrm{total}}$ (e.g.\ $K_{\mathrm{total}} = 960$ split into thirds).
  \item For each profile, time indices are sorted by profile value and split into $B$ equal-population bins; \textbf{anchors} (indices of min and max) are always included. From each bin, a random subset of indices is drawn without replacement until the quota $K_{\mathrm{load}}$, $K_{\mathrm{pv}}$, or $K_{\mathrm{net}}$ is filled. The union of the three sets is taken (duplicates removed); if size $\ne K_{\mathrm{total}}$, it is trimmed or augmented randomly.
\end{itemize}
Result: a set of time indices $\mathcal{T}_s \subseteq \{0,\ldots,T-1\}$ with $|\mathcal{T}_s| \le K_{\mathrm{total}}$ per scenario $s$.

\section{Device-level application at time $t$}

For each selected $t$, load and PV are applied to the circuit as follows.

\subsection{Load}

For each load device $d$ with share $\alpha_d^P$, $\alpha_d^Q$ (from \texttt{DEVICE\_P\_SHARE}, \texttt{DEVICE\_Q\_SHARE}; $\sum_d \alpha_d^P = \sum_d \alpha_d^Q = 1$):
\begin{equation}
  P_d(t) = P_{\mathrm{load}}(t)\, \alpha_d^P\, (1 + \varepsilon_d^P), \qquad
  Q_d(t) = Q_{\mathrm{load}}(t)\, \alpha_d^Q\, (1 + \varepsilon_d^Q),
\end{equation}
where $\varepsilon_d^P, \varepsilon_d^Q \sim \mathcal{N}(0, \sigma_{\mathrm{load}})$ i.i.d., and the factor $(1+\varepsilon)$ is clipped to $\ge 0$. These $P_d(t)$, $Q_d(t)$ are written to OpenDSS. Node-wise load (and later PV) is obtained by aggregating device-level values to bus-phase nodes as follows.

\paragraph{Device-to-node aggregation.} Each load (and PV) device is mapped to one or more $(bus, phase)$ nodes with a weight $w$ summing to 1 per device. The implementation uses \texttt{DEVICE\_TO\_BUSPH} (and \texttt{PV\_TO\_BUSPH}): a list of $(bus, phase, w)$ per device. Three-phase loads assign one third of the device's $P_d$, $Q_d$ to each of the three phase nodes at that bus ($w = 1/3$ per phase). Single-phase loads assign 100\% to the corresponding $(bus, phase)$ node ($w = 1$). Line or transformer loads that sit between buses may split 0.5/0.5 across the two bus-phase endpoints. Then
\begin{equation}
  P_{\mathrm{load},i}(t) = \sum_{d \colon (bus_i, ph_i) \in d} P_d(t)\, w_{d,i}, \qquad
  Q_{\mathrm{load},i}(t) = \sum_{d \colon (bus_i, ph_i) \in d} Q_d(t)\, w_{d,i},
\end{equation}
where $i$ indexes the node $(bus_i, ph_i)$ and $w_{d,i}$ is the weight of device $d$ for that node. The same aggregation is used for PV to get $P_{\mathrm{pv},i}(t)$. This defines the node-level quantities that enter the injection features below; it is implemented in the dataset code when building \texttt{busphP\_load}, \texttt{busphQ\_load}, and the PV equivalents.

\subsection{PV}

PV devices are assigned a fraction of total $P_{\mathrm{pv}}$ (e.g.\ \texttt{PV\_PMMP\_SHARE}); at time $t$ the Pmpp of each PV system is set so that its effective output (after the irradiance shape in OpenDSS) is $P_{\mathrm{pv}}(t)\, \beta_d\, (1+\varepsilon_d^{\mathrm{pv}})$, with $\varepsilon_d^{\mathrm{pv}} \sim \mathcal{N}(0, \sigma_{\mathrm{pv}})$ and factor $\ge 0$. The resulting per-node PV injection $P_{\mathrm{pv},i}(t)$ is obtained from \texttt{PV\_TO\_BUSPH}.

\subsection{Capacitors}

Fixed capacitor Q at specified buses (e.g.\ 844, 848) is read from \texttt{CAP\_Q\_KVAR} and added to reactive injection at those bus-phase nodes.

\section{Node-level injection (features)}

After solving the AC power flow, the \textbf{injection features} stored per node $i$ are:
\begin{equation}
  P_{\mathrm{inj},i}(t) =
  \begin{cases}
    P_{\mathrm{grid}}(t)\,/\,3 & \text{(source bus, 3 phases)} \\
    P_{\mathrm{pv},i}(t) - P_{\mathrm{load},i}(t) & \text{otherwise},
  \end{cases}
\end{equation}
\begin{equation}
  Q_{\mathrm{inj},i}(t) =
  \begin{cases}
    Q_{\mathrm{grid}}(t)\,/\,3 & \text{(source bus)} \\
    -Q_{\mathrm{load},i}(t) + Q_{\mathrm{cap},i} & \text{otherwise},
  \end{cases}
\end{equation}
where $P_{\mathrm{grid}}(t)$, $Q_{\mathrm{grid}}(t)$ are the total active and reactive power from the \textbf{upstream grid} (positive = power supplied into the feeder; obtained from OpenDSS \texttt{TotalPower} and negated). At the source bus, the injection is thus the \textbf{power from upstream}, split equally across the three phases ($\div 3$). At all other nodes, injection is net local (PV minus load, and for $Q$: minus load plus capacitor). $Q_{\mathrm{cap},i}$ is the fixed capacitor Q at that bus (if any), zero otherwise.

\section{Solve, targets, and rejection}

\begin{enumerate}[noitemsep]
  \item Set simulation time to step $t$, apply $P_d(t)$, $Q_d(t)$ and PV as above.
  \item Run one AC power-flow solve.
  \item If not converged: discard this $(s,t)$; no sample.
  \item Read voltage magnitude $|V_i|$ (pu) and angle $\angle V_i$ (deg) at all nodes. If any $|V_i| \notin [V_{\min}, V_{\max}]$ (e.g.\ $[0.5, 1.5]$ pu) or non-finite: discard.
  \item Otherwise: append one sample; store per-node rows with features \texttt{p\_inj\_kw}, \texttt{q\_inj\_kvar} and targets \texttt{vmag\_pu}, \texttt{vang\_deg}.
\end{enumerate}

\section{Outputs}

\begin{itemize}[noitemsep]
  \item \textbf{gnn\_node\_index\_master.csv}: node name and node index for all $N$ nodes.
  \item \textbf{gnn\_edges\_phase\_static.csv}: static phase-level graph (lines and transformers).
  \item \textbf{gnn\_sample\_meta.csv}: one row per kept sample: \texttt{sample\_id}, \texttt{scenario\_id}, \texttt{t\_index}, \texttt{t\_minutes}, totals and profile values at $t$.
  \item \textbf{gnn\_node\_features\_and\_targets.csv}: one row per (sample, node): \texttt{sample\_id}, \texttt{node}, \texttt{node\_idx}, \texttt{p\_inj\_kw}, \texttt{q\_inj\_kvar}, \texttt{vmag\_pu}, \texttt{vang\_deg}.
\end{itemize}
The GNN is trained to predict \texttt{vmag\_pu} (and optionally \texttt{vang\_deg}) from the injection features and the static graph.

\end{document}
